% !TeX root = ../泛函分析笔记.tex
% 由主文档通过 \input 载入；请编译主文档。
\setcounter{chapter}{3} % 保留原书章号。
\chapter{部分习题解答与证明方法回顾}
\label{app:fa-solutions}

本附录选取能够代表主要证明方法的习题，给出完整论证。阅读时宜先独立确定所需估计，再与解答比较；只记最后的结论，难以掌握这些方法在新问题中的用法。其余习题可沿相邻例题与定理的证明尝试。

\section{由局部估计控制整个序列}
第\ref{chap:fa-metric}章习题\ref{exer:fa-cauchy}要求证明 Cauchy 序列有界，并说明一个收敛子列足以决定整个 Cauchy 序列的极限。取 $\varepsilon=1$，存在 $N$ 使 $m,n\geq N$ 时 $d(x_m,x_n)<1$。固定 $m=N$，后面所有项都落在 $B(x_N,1)$ 中；前面只有有限项，取
\[
 R=1+\max_{1\leq j\leq N}d(x_j,x_N),
\]
即可把整个序列包含在 $B(x_N,R)$ 中。这一步使用的是“尾部统一受控，加上有限个例外”的结构。

若 $x_{n_k}\to x$，给定 $\varepsilon>0$，先由 Cauchy 条件选择 $N$，使后部两项距离小于 $\varepsilon/2$，再选 $k$ 使 $n_k\geq N$ 且 $d(x_{n_k},x)<\varepsilon/2$。对所有 $n\geq N$，三角不等式给出 $d(x_n,x)<\varepsilon$。注意这里不是给每个 $n$ 临时寻找一个不同的近点，而是用同一个子列元素作为整个尾部的比较点。这个论证经常用于完备性证明：先为一个精心选择的子列找到极限，再回到原序列。

\section{数列范数的包含关系}
第\ref{chap:fa-banach}章习题\ref{exer:fa-lp-inclusion}可以先通过归一化消除大小因素。若 $x=0$，结论直接成立；否则令 $z=x/\|x\|_p$，则 $\sum_j|z_j|^p=1$，所以每个 $|z_j|\leq1$。当 $q<\infty$ 时，$|z_j|^q\leq|z_j|^p$，故 $\sum_j|z_j|^q\leq1$，从而 $\|x\|_q\leq\|x\|_p$。当 $q=\infty$ 时，逐坐标估计直接给出同一结论。

当 $p<q<\infty$，选择 $1/q<a<1/p$，则 $(j^{-a})\in\ell^q$ 而不在 $\ell^p$，说明包含通常严格。$q=\infty$ 时，常数列已经给出严格性。这里的证明依赖计数测度的每个单点都有固定正质量，因此一个坐标不能在 $\ell^p$ 范数有界时任意变大；有限区间上的 $L^p$ 函数却可在很小的集合上取很大的值，这正是两类包含方向不同的原因。

\section{伴随与正交补的计算}
第\ref{chap:fa-hilbert}章习题\ref{exer:fa-shift-adjoint}中，Cauchy--Schwarz 保证下列级数绝对收敛，因此移换指标合法：
\[
 \langle Rx,y\rangle=\sum_{j\geq2}x_{j-1}\overline{y_j}
 =\sum_{k\geq1}x_k\overline{y_{k+1}}=\langle x,Ly\rangle.
\]
由伴随唯一性，$R^*=L$、$L^*=R$。又 $LR=I$，而 $RLx=x-x_1e_1$，所以右移保持范数，却不满射，且 $RR^*\ne I$，不是酉算子。等距算子只要求 $T^*T=I$；酉算子还要求能够从每个输出倒推输入，也就是在此基础上满射。

习题\ref{exer:fa-double-perp}则适合使用正交分解。首先 $M^\perp=(\overline M)^\perp$，因为内积连续；每个 $m\in\overline M$ 与 $M^\perp$ 正交，所以 $\overline M\subset M^{\perp\perp}$。反向对 $x\in M^{\perp\perp}$，分解 $x=m+z$，其中 $m\in\overline M$、$z\in M^\perp$。由于 $x$ 与 $z$ 正交，而 $m$ 也与 $z$ 正交，便有 $0=\langle x,z\rangle=\|z\|^2$，所以 $z=0$、$x=m\in\overline M$。两次正交补会自动补齐极限点，因此结果是闭包，未必是原来的非闭子空间。

\section{Baire 方法与稠密集上的验证}
第\ref{chap:fa-baire}章习题\ref{exer:fa-hamel}中，反设 $e_1,e_2,\ldots$ 为无限维 Banach 空间 $X$ 的可数 Hamel 基。令 $M_n=\operatorname{span}\{e_1,\ldots,e_n\}$，每个 $M_n$ 有限维，因而闭。真线性子空间内部必为空：若 $B(x_0,r)\subset M_n$，相减得到 $B(0,r)\subset M_n$，再由数乘可把每个向量缩入该球，迫使 $M_n=X$，矛盾。于是每个 $M_n$ 无处稠密。但 Hamel 展开仅有有限项，所以 $X=\bigcup_nM_n$，违反 Baire 定理。可数标准正交基则允许无穷级数，不给出这个可数并等式，因此不会产生矛盾。

习题\ref{exer:fa-dense-convergence}设 $D\subset X$ 稠密，$Y$ 完备，且 $\|T_n\|\leq M$。给定 $x\in X$ 与 $z\in D$，
\[
 \|T_nx-T_mx\|\leq2M\|x-z\|+\|T_nz-T_mz\|.
\]
若 $M=0$，所有算子都为零；否则先取 $z$ 使第一项小于 $\varepsilon/2$，再由 $D$ 上的收敛选择 $n,m$ 使第二项小于 $\varepsilon/2$。于是 $(T_nx)$ 在 $Y$ 中 Cauchy，故收敛。逐点极限 $T$ 线性且 $\|Tx\|\leq M\|x\|$。这项结论只需要值域完备；若统一算子界已经给出，定义域本身无需完备。统一估计负责把稠密集中的结论传到全空间，完备性负责接住输出的极限。

\section{弱紧性怎样保证最值}
第\ref{chap:fa-reflexive}章习题\ref{exer:fa-norm-attainment}中，若 $f=0$，任取 $x=0$ 即可。若 $f\ne0$，按上确界定义选 $x_n\in B_X$，使 $|f(x_n)|\to\|f\|$。自反性给出弱收敛子列 $x_{n_k}\rightharpoonup x$；范数弱下半连续，故 $\|x\|\leq1$。连续线性泛函按弱收敛的定义满足 $f(x_{n_k})\to f(x)$，取绝对值得 $|f(x)|=\|f\|$。由于 $|f(x)|\leq\|f\|\|x\|$，进一步得到 $\|x\|=1$。

这个论证与紧集上连续函数取得最大值的证明相同，只是把范数子列换成弱收敛子列。不能把结论推广到单位球上的任意范数连续函数，因为范数连续不保证弱连续；例如无限维 Hilbert 空间中的范数函数在 $e_n\rightharpoonup0$ 时不保持函数值极限。

\section{Galerkin 收敛中嵌套性的作用}
第\ref{chap:fa-variational}章习题\ref{exer:fa-galerkin}中，已知 $H_n\subset H_{n+1}$ 且 $\bigcup_nH_n$ 稠密。给定 $\varepsilon>0$，稠密性保证存在某个 $N$ 与 $w\in H_N$，使 $\|u-w\|<\varepsilon$；嵌套性保证同一个 $w$ 属于所有 $n\geq N$ 的 $H_n$。因此这些 $n$ 都满足
\[
 \inf_{w_n\in H_n}\|u-w_n\|\leq\|u-w\|<\varepsilon.
\]
Céa 估计随即给出 $u_n\to u$。这里先找到一个好近似，再保证它在以后所有逼近空间中仍可使用，正是嵌套性的作用。

仅有并集稠密却未必足够。在 $\ell^2$ 中，可令 $H_{2n}=\operatorname{span}\{e_1,\ldots,e_n\}$，$H_{2n-1}=\operatorname{span}\{e_1\}$。它们的并集稠密，但对于 $u=e_2$，奇数步的最佳逼近误差始终为一。若取 $a(u,v)=\langle u,v\rangle$ 的实 Hilbert 空间问题，Galerkin 解就是正交投影，因此这一序列的离散解也不收敛。非嵌套方法并非不可用，但必须直接验证对整个指标序列成立的逼近条件。

